\begingroup
\renewcommand{\arraystretch}{1.45}
\clearpage\fancyhead[L]{\small REMARKS.}\pdfbookmark[0]{Remarks on Interpolation}{interpolation-bemerkungen}\SourceLeaf{337}
\begin{center}\textbf{REMARKS.}\end{center}
\begin{center}\rule{3em}{0.3pt}\end{center}

The present treatise on interpolation is a reworking of an investigation of which an earlier draft seems to have been begun in October 1805. It contains the methods that have often been applied in calculating the expressions for the changes in the elements of a planetary orbit produced by perturbing forces. Alongside this determination of the perturbations stands another procedure, which directly gives the values of the orbital elements for individual points in time and which, as follows from the final formulas recorded in the manuscript papers found among the posthumous remains, rests on the following principles.

\Needspace{28\baselineskip}
If the values of the argument for which the values of the function are given form an arithmetic progression, and one has introduced an argument $x$ which, at those values, advances like the integers, denotes the given values by $fx$, introduces their successive difference series:
\[
\begin{aligned}
&f(x+1)-fx=f^{1}(x+\frac12),\\[6pt]
&f^{1}(x+\frac12)-f^{1}(x-\frac12)=f^{2}x,\\[6pt]
&f^{2}(x+1)-f^{2}x=f^{3}(x+\frac12)\ \mbox{etc.}
\end{aligned}
\]
then their sum series beginning with arbitrary initial terms,
\hfill\break{}
$f^{-1}(x+\nobreak{}\tfrac12),\allowbreak{}\ f^{-2}x,\allowbreak{}\ f^{-3}(x+\nobreak{}\tfrac12)$ etc.,
for which
\[
\begin{aligned}
&fx=f^{-1}(x+\frac12)-f^{-1}(x-\frac12),\\[6pt]
&f^{-1}(x+\frac12)=f^{-2}(x+1)-f^{-2}x\ \mbox{etc.}
\end{aligned}
\]
and finally the mean values of successive terms of the series thus formed, namely
\[
\begin{aligned}
&\mathinner{.\,.\,.} f^{-1}x=\frac12f^{-1}(x+\frac12)+\frac12f^{-1}(x-\frac12),\\[6pt]
&f(x+\frac12)=\frac12f(x+1)+\frac12fx,\\[6pt]
&f^{1}x=\frac12f^{1}(x+\frac12)+\frac12f^{1}(x-\frac12)\mathinner{.\,.\,.}
\end{aligned}
\]
thus, arranged for a clear overview, the following system of given values and values initially calculated from them
\SourceLeaf{338}

\begin{landscape}\thispagestyle{empty}\begingroup\fontsize{9}{12}\selectfont\setlength{\arraycolsep}{3pt}\renewcommand{\arraystretch}{1.3}
\[
\begin{array}{ccc|cc||cc||cc|ccc}
 & \mbox{M. V.} & \mbox{II. sum.} & \mbox{M. V.} & \mbox{I. sum} & \mbox{given v.} & \mbox{M. V.} & \mbox{I. diff.} & \mbox{M. V.} & \mbox{II. diff.} & \mbox{M. V.} & \\
 & & \cdot & \cdot & & \cdot & & & \cdot & \cdot & & \\
 & \cdot & & & \cdot & & \cdot & \cdot & & & \cdot & \cdot \\
 & & f^{-2}(x-1) & f^{-1}(x-1) & & f(x-1) & & & f^{1}(x-1) & f^{2}(x-1) & & \\
\cdot & f^{-2}(x-\tfrac12) & & & f^{-1}(x-\tfrac12) & & f(x-\tfrac12) & f^{1}(x-\tfrac12) & & & f^{2}(x-\tfrac12) & \cdot \\
 & & f^{-2}x & f^{-1}x & & fx & & & f^{1}x & f^{2}x & & \\
\cdot & f^{-2}(x+\tfrac12) & & & f^{-1}(x+\tfrac12) & & f(x+\tfrac12) & f^{1}(x+\tfrac12) & & & f^{2}(x+\tfrac12) & \cdot \\
 & & f^{-2}(x+1) & f^{-1}(x+1) & & f(x+1) & & & f^{1}(x+1) & f^{2}(x+1) & & \\
\cdot & \cdot & & & \cdot & & \cdot & \cdot & & & \cdot & \cdot \\
 & & \cdot & \cdot & & \cdot & & & \cdot & \cdot & &
\end{array}
\]
\endgroup\end{landscape}

\Needspace{30\baselineskip}
then Article 4, if one $1^{\mathrm{st}}$ substitutes for $t,\allowbreak{}\ a,\allowbreak{}\ b,\allowbreak{}\ c,\allowbreak{}\ d,\allowbreak{}\ e\ldotp\ldotp$
respectively $x+t,\allowbreak{}\ x,\allowbreak{}\ x+1,\allowbreak{}\ x-1,\allowbreak{}\ x+2,\allowbreak{}\ x-2$, $2^{\mathrm{nd}}$ for
$t,\allowbreak{}\ a,\allowbreak{}\ b,\allowbreak{}\ c,\allowbreak{}\ d,\allowbreak{}\ e\ldotp\ldotp$ respectively $x+t,\allowbreak{}\ x,\allowbreak{}\ x-1,\allowbreak{}\ x+1,\allowbreak{}\ x-2,\allowbreak{}\ x+2\mathinner{.\,\allowbreak{}.\,\allowbreak{}.}$,
$3^{\mathrm{rd}}$ forms half the sum of the corresponding sides of the two equations thus obtained,
$4^{\mathrm{th}}$ forms half the sum from the same equations, after first substituting in the second $x+1$ and $t-1$ respectively for $x$ and $t$,
and $5^{\mathrm{th}}$ puts $t=\tfrac12$ in the equation last obtained, gives
the following equations for the values $\varphi(x+\nobreak{}t)$ and
$\varphi(x+\nobreak{}\tfrac12)$ to be determined by interpolation, of the function $\varphi x$ which, at every integer value
of $x$, becomes equal to the given quantity $fx$:
\[
\begin{aligned}
&\varphi(x+t)={}fx+t\cdot f^{1}(x+\frac12)+\frac{t\cdot t-1}{1\cdot2}f^{2}x
\\[6pt]
&\qquad{}+\frac{t\cdot t-1\cdot t+1}{1\cdot2\cdot3}f^{3}(x+\frac12)+\ldotp\ldotp\\[6pt]
&+\frac{\Pi(t+n-1)}{\Pi 2n\cdot\Pi(t-n-1)}f^{2n}x
\\[6pt]
&\qquad{}+\frac{\Pi(t+n)}{\Pi(2n+1)\cdot\Pi(t-n-1)}f^{2n+1}(x+\frac12)\ldotp\ldotp
\end{aligned}
\]
\Needspace{18\baselineskip}
\[
\begin{aligned}
&\varphi(x+t)={}fx+t\cdot f^{1}(x-\frac12)+\frac{t\cdot t+1}{1\cdot2}f^{2}x
\\[6pt]
&\qquad{}+\frac{t\cdot t+1\cdot t-1}{1\cdot2\cdot3}f^{3}(x-\frac12)+\ldotp\ldotp\\[6pt]
&+\frac{\Pi(t+n)}{\Pi 2n\cdot\Pi(t-n)}f^{2n}x
\\[6pt]
&\qquad{}+\frac{\Pi(t+n)}{\Pi(2n+1)\cdot\Pi(t-n-1)}f^{2n+1}(x-\frac12)\mathinner{.\,.\,.}
\end{aligned}
\]
\[
\begin{aligned}
&\varphi(x+t)={}fx+t\cdot f^{1}x+\frac{t\cdot t}{1\cdot2}f^{2}x
+\frac{t\cdot t-1\cdot t+1}{1\cdot2\cdot3}f^{3}x+\ldotp\ldotp\\[6pt]
&+\frac{\Pi(t+n-1)}{\Pi 2n\cdot\Pi(t-n)}t f^{2n}x
+\frac{\Pi(t+n)}{\Pi(2n+1)\Pi(t-n-1)}f^{2n+1}x\\[6pt]
&\qquad{}+\mathinner{.\,.\,.}
\end{aligned}
\]
\Needspace{22\baselineskip}
\[
\begin{aligned}
&\varphi(x+t)={}f(x+\frac12)+(t-\frac12)f^{1}(x+\frac12)
+\frac{t\cdot t-1}{1\cdot2}f^{2}(x+\frac12)
\\[6pt]
&\qquad{}+\frac{t\cdot t-1}{1\cdot2}(t-\frac12)f^{3}(x+\frac12)+\ldotp\ldotp\\[6pt]
&+\frac{\Pi(t+n-1)}{\Pi 2n\cdot\Pi(t-n-1)}f^{2n}(x+\frac12)
\\[6pt]
&\qquad{}+\frac{\Pi(t+n-1)}{\Pi(2n+1)\cdot\Pi(t-n-1)}(t-\frac12)f^{2n+1}(x+\frac12)+\ldotp\ldotp
\end{aligned}
\]
\[
\begin{aligned}
&\varphi(x+\frac12)={}f(x+\frac12)-\frac18f^{2}(x+\frac12)
+\frac{3}{128}f^{4}(x+\frac12)\\[6pt]
&\qquad{}-\frac{5}{1024}f^{6}(x+\frac12)
+\frac{35}{32768}f^{8}(x+\frac12)\mathinner{.\,.\,.}\\[6pt]
&+(-1)^n\frac{\Pi 2n}{\Pi n\cdot\Pi n}\,2^{-4n}\cdot f^{2n}(x+\frac12)+\mathinner{.\,.\,.}
\end{aligned}
\]
\Needspace{14\baselineskip}
The coefficients of $f^m$ in the third, fourth and fifth equations are equal to the coefficients of
$(2i\sin\omega)^m$ in the series expansions respectively of
$\cos 2t\omega+i\dfrac{\sin 2t\omega}{\cos\omega},\allowbreak{}\allowbreak \quad
\dfrac{\cos(2t-\nobreak{}1)\omega}{\cos\omega}+i\sin(2t-\nobreak{}1)\omega$
and $\dfrac{1}{\cos\omega}$ in powers of the quantity $2i\sin\omega$.
\Needspace{28\baselineskip}
\par\SourceLeaf{339}
If one now sets
\[
\begin{aligned}
&\Phi(x+\frac12)={}f^{-1}(x+\frac12)+\frac{1}{24}f^{1}(x+\frac12)
-\frac{17}{5760}f^{3}(x+\frac12)\\[6pt]
&\qquad{}+\frac{367}{967680}f^{5}(x+\frac12)
-\frac{27859}{164486400}f^{7}(x+\frac12)\\[6pt]
&+\frac{1295803}{122624409600}f^{9}(x+\frac12)-\ldotp\ldotp
\end{aligned}
\]
\[
\begin{aligned}
&\Phi x={}f^{-1}x-\frac{1}{12}f^{1}x+\frac{11}{720}f^{3}x
-\frac{191}{60480}f^{5}x+\frac{2497}{3628800}f^{7}x\\[6pt]
&-\frac{14797}{95800320}f^{9}x+\frac{92427157}{2615348736000}f^{11}x-\mathinner{.\,.\,.}
\end{aligned}
\]
\[
\begin{aligned}
&\Psi x={}f^{-2}x+\frac{1}{12}fx-\frac{1}{240}f^{2}x
+\frac{31}{60480}f^{4}x-\frac{289}{3628800}f^{6}x\\[6pt]
&+\frac{317}{22809600}f^{8}x-\ldotp\ldotp
\end{aligned}
\]
\[
\begin{aligned}
&\Psi(x+\frac12)=f^{-2}(x+\frac12)-\frac{1}{24}f(x+\frac12)
+\frac{17}{1920}f^{2}(x+\frac12)\\[6pt]
&\qquad{}-\frac{367}{193536}f^{4}(x+\frac12)+\ldotp-\ldotp\ldotp
\end{aligned}
\]
\Needspace{30\baselineskip}
taking as the coefficients of $f^m$ in these four equations the coefficients that, in the
expansions respectively of
$\dfrac{1}{2i\omega},\allowbreak{}\allowbreak \ \dfrac{1}{2i\omega\cos\omega},\allowbreak{}\allowbreak \ -\dfrac{1}{4\omega\omega},\allowbreak{}\allowbreak \ -\dfrac{1}{4\omega\omega\cos\omega}$,
still arise in powers of $2i\sin\omega$,
then $\Phi(x+\nobreak{}\tfrac12),\allowbreak{}\ \Phi x,\allowbreak{}\ \Psi x,\allowbreak{}\ \Psi(x+\nobreak{}\tfrac12)$ are the particular values of functions
$\Phi(x+\nobreak{}t),\allowbreak{}\ \Psi(x+\nobreak{}t)$
for which the equations
\[
\begin{aligned}
&\Phi(x+t)-\Phi(x+t_0)=\int_{t_0}^{t}\varphi(x+t)\,dt),\\[6pt]
&\Psi(x+t)-\Psi(x+t_1)=\int_{t_1}^{t}\Phi(x+t)\,dt
\end{aligned}
\]
hold; they therefore become equal to definite integrals of $\varphi(x+\nobreak{}t)\,\allowbreak{}dt$ and
$\varphi(x+\nobreak{}t)\cdot dt\cdot dt$ if
the initial values of the sum series
\hfill\break{}
$f^{-1}(x-\nobreak{}\tfrac12),\allowbreak{}\ f^{-1}(x+\nobreak{}\tfrac12)\ldotp\ldotp,\allowbreak{}\ f^{-2}x,\allowbreak{}\ f^{-2}(x+\nobreak{}1)\ldotp\ldotp$
are suitably chosen.

The derivation of the series for $\Phi$ and $\Psi$ from the third and fourth of the above equations for
$\varphi(x+\nobreak{}t)$ is obtained directly if the integrations are carried out so that the expressions for
\hfill\break{}
$\Phi(x+\nobreak{}\tfrac12)-\Phi(x-\nobreak{}\tfrac12),\allowbreak{}\quad \Phi(x+\nobreak{}1)-\Phi x,\allowbreak{}\quad
\Psi(x+\nobreak{}1)-2\Psi x+\Psi(x-\nobreak{}1)$ and
$\Psi(x+\nobreak{}\tfrac32)-2\Psi(x+\nobreak{}\tfrac12)+\Psi(x-\nobreak{}\tfrac12)$
result first.
\begin{flushright}\textsc{Schering.}\end{flushright}
\begin{center}\rule{9em}{0.3pt}\end{center}
\clearpage
\endgroup
